\begin{center}\begin{tikzpicture}[font=\small,>=stealth,line join=round]
\draw[gray!40,rounded corners=3pt] (0,0) rectangle (14.2,9.2);\node at (7.1,8.6) {高度相同：一直增大，两端都只能靠近};
\node[blue] at (7.1,7.65) {$\displaystyle h(x)=\sqrt{x^2+4}-\sqrt{(x-3)^2+4}$};\begin{scope}[shift={(4.1,4.0)},x=0.3130434782608696cm,y=0.75cm]\draw[->,gray] (-10,0)--(13.5,0) node[right]{$x$};\draw[->,gray] (0,-3.6)--(0,3.8000000000000003) node[left]{$y$};\draw[gray,dashed] (-10,-3)--(13,-3);\node[left=3pt,fill=white,inner sep=1pt] at (0,-3) {$-3$};\draw[gray,dashed] (-10,3)--(13,3);\node[left=3pt,fill=white,inner sep=1pt] at (0,3) {$3$};\draw[blue,very thick,domain=-10:13,samples=240] plot(\x,{sqrt(\x*\x+4)-sqrt((\x-3)^2+4)});\end{scope}\node[anchor=west,align=left] at (9,4.5) {$a=b=2,\ m=0,n=3$\\[12pt]对每个实数 $x$\\[10pt]都有 $-3<h(x)<3$\\[12pt]$h(3/2)=0$\\[12pt]完整值域：$(-3,3)$};
\node[text width=12.8cm,align=center] at (7.1,.65) {两条虚线均不属于函数图像；同高度且不同中心时没有最大、最小值。};
\end{tikzpicture}\end{center}
